\section{Introduction}\label{sec:intro}

Systems that answer questions over PDFs first extract the text, then cut it
into chunks that are indexed and retrieved. Both steps can lose information:
a parser can drop or garble text, and a chunker can split the evidence for an
answer across pieces. Because documents differ in layout and structure, a
natural idea is to choose the parser and chunker for each document, and recent
work proposes doing so from intrinsic quality signals
\cite{demourajunior2026adaptive}.

Before building such a chooser, we ask a narrower question: how much
retrieval quality is lost by using one fixed parse-and-chunk setup for every
PDF, and does a choice per document, made by a cross-fitted oracle from a part
of each document's own questions, recover any of it? The oracle selects on few
questions, so its gap is not a bound on every chooser. We study this together with
four related questions:
\begin{itemize}
\item \textbf{RQ1:} Does choosing a setup per document improve retrieval over
the best single setup?
\item \textbf{RQ2:} At a fixed token budget, how does chunk size relate to
recall?
\item \textbf{RQ3:} Does the best chunk size change when retrieval is measured
with a fixed number of chunks instead of a fixed token budget?
\item \textbf{RQ4:} Do differences in answer quality follow differences in
retrieval?
\item \textbf{RQ5:} Do tuned parser settings and chunk sizes beat their
defaults (a library default for the parsers and LangChain, our own starting
size for the other chunkers)?
\end{itemize}

We compare setups by the text they let a reader see, not by how many chunks
they return: every setup reads its ranked chunks up to the same token budget.
The RQ1 design, the scorer and the RQ1 decision rule were written into dated
specifications before the first held-out set was scored; the designs of the
other parts came later and are dated in Section~\ref{sec:method}. The NO-GO
decision is repeated on a second held-out set with a different source of gold
evidence.
Where the study departed from its plan, we report it.

Our contributions are:
\begin{enumerate}
\item A pre-registered test of per-document selection on two sealed held-out
sets, with an oracle that is cross-fitted so that it does not gain from
selection luck. It finds no evidence of a gain over a single small-chunk setup
under a fixed token budget (Section~\ref{sec:rq1}).
\item A token-budgeted comparison of chunk sizes, showing that small chunks
tend to retrieve best at a budget of 750 or 1,250 tokens (not established at
2,000) and, on tuning data, that a fixed number of chunks reverses this ranking
(Sections~\ref{sec:rq2} and~\ref{sec:rq3}).
\item A comparison of tuning against defaults: LangChain's default chunk size
costs recall (post-hoc, on held-out set~1), while parser settings show no
detectable effect, and for scans what matters is detecting them
(Section~\ref{sec:rq5}).
\item A full record of deviations, reuse of held-out data and post-hoc
analyses, with the code, the checksummed results and a tool with fixed defaults.
\end{enumerate}
